\section{群作用与有限群结构}
\label{chap:finite-group-structure}

群的乘法、陪集和商群描述了群本身的结构。对称性的物理内容还需要说明群元对别的对象做了什么。一个转动可以移动顶点，一个置换可以移动标签；轨道收集所有可达对象，稳定子则记录哪些操作把选定对象留在原处。

把群作用在自身上并取共轭，就得到共轭类和类方程。Cauchy 与 Sylow 定理正是在这套计数语言下限制有限群的素数阶子群。直积和半直积放在最后，用来反过来组装群；二面体群会显示“一个子群怎样扭转另一个正规子群”。

\subsection{变换群与群作用}
\label{sec:group-actions}

到这里为止，群似乎一直是一个“自己和自己运算”的对象。但物理上我们真正关心的通常不是群元本身，而是它对别的对象做了什么：转动怎样移动空间中的点，点群怎样变换原子轨道，置换怎样交换粒子标签。把“群对另一个集合的作用”抽象出来，就得到群作用。

群作用最关键的两条规则与表示的同态条件非常相似：单位元必须什么都不改变，复合群元 $gh$ 的作用必须等于先做 $h$ 再做 $g$。区别在于，这里的被作用对象 $X$ 暂时只要求是集合，不必有向量加法。下一章把 $X$ 换成线性空间并要求作用是线性的，就会得到群表示。

“变换群”的核心可以用群作用统一表达。这样做的好处是：变换对象不必是几何点，也可以是晶格、波函数、子空间、哈密顿量或一组基矢；同一套轨道--稳定子语言都适用。

\begin{definitionplain}[群作用]
设 $G$ 为群，$X$ 为非空集合。$G$ 在 $X$ 上的一个左作用是映射
\[
G\times X\longrightarrow X,
\qquad
(g,x)\longmapsto g\cdot x,
\]
满足
\[
e\cdot x=x,
\qquad
(g_1g_2)\cdot x=g_1\cdot(g_2\cdot x).
\label{eq:group-action}
\]
等价地，一个群作用就是一个同态
\[
\rho:G\longrightarrow\operatorname{Sym}(X).
\]
\label{def:group-action}
\end{definitionplain}

\begin{definitionplain}[忠实作用]
若作用对应的同态 $\rho:G\to\operatorname{Sym}(X)$ 为单射，则称该作用为忠实作用。等价地，只有单位元在 $X$ 上处处作用为恒等变换。
\label{def:faithful-action}
\end{definitionplain}

\begin{definitionplain}[轨道与稳定子]
给定 $x\in X$，其轨道定义为
\[
\operatorname{Orb}_G(x)=\{g\cdot x\mid g\in G\}.
\]
其稳定子（亦称迷向子群）定义为
\[
\operatorname{Stab}_G(x)=\{g\in G\mid g\cdot x=x\}.
\]
\label{def:orbit-stabilizer}
\end{definitionplain}

\begin{proposition}[稳定子是子群]
对任意 $x\in X$，
\[
\operatorname{Stab}_G(x)\leqslant G.
\]
\label{prop:stabilizer-subgroup}

\end{proposition}

\begin{proof}
若 $a,b\in\operatorname{Stab}_G(x)$，则
\[
(ab^{-1})\cdot x
=a\cdot(b^{-1}\cdot x).
\]
由 $b\cdot x=x$，两边作用 $b^{-1}$ 得 $b^{-1}\cdot x=x$，于是
\[
(ab^{-1})\cdot x=a\cdot x=x.
\]
故 $ab^{-1}\in\operatorname{Stab}_G(x)$。
\end{proof}

\begin{theorem}[轨道--稳定子定理]
对任意群作用与 $x\in X$，存在自然双射
\[
\{g\operatorname{Stab}_G(x)\mid g\in G\}
\longleftrightarrow
\operatorname{Orb}_G(x),
\qquad
g\operatorname{Stab}_G(x)\longmapsto g\cdot x.
\label{eq:orbit-stabilizer-bijection}
\]
左边只是稳定子的左陪集集合。稳定子未必是正规子群，所以这里不需要、也通常不存在商群结构。
若 $G$ 有限，则
\[
\lvert\operatorname{Orb}_G(x)\rvert
=[G:\operatorname{Stab}_G(x)]
=\frac{\lvert G\rvert}{\lvert\operatorname{Stab}_G(x)\rvert}.
\label{eq:orbit-stabilizer}
\]
\label{thm:orbit-stabilizer}

\end{theorem}

\begin{proof}
若 $g\operatorname{Stab}_G(x)=h\operatorname{Stab}_G(x)$，则 $h^{-1}g\in\operatorname{Stab}_G(x)$，故
\[
h^{-1}g\cdot x=x
\Longrightarrow
g\cdot x=h\cdot x,
\]
所以映射良定义。任意轨道元素按定义形如 $g\cdot x$，故满射。反之若 $g\cdot x=h\cdot x$，则
\[
h^{-1}g\cdot x=x,
\]
即 $h^{-1}g\in\operatorname{Stab}_G(x)$，因此两陪集相同，故单射。有限情形再取集合基数即得\eqnrefs{eq:orbit-stabilizer}。
\end{proof}

证明中的关键映射是
\[
\{g\operatorname{Stab}_G(x)\mid g\in G\}\longrightarrow \operatorname{Orb}_G(x),
\qquad g\operatorname{Stab}_G(x)\longmapsto g\cdot x.
\]
这里再强调一次：斜线记号若偶尔写成 $G/\operatorname{Stab}_G(x)$，只应读作“稳定子的左陪集集合”，除非已经额外证明 $\operatorname{Stab}_G(x)\mathrel{\trianglelefteq} G$，否则不能把它当作商群。
为什么要用陪集而不是单独的 $g$？因为不同群元可能把 $x$ 送到同一点。若 $g_1\cdot x=g_2\cdot x$，则
\[
g_2^{-1}g_1\cdot x=x,
\]
所以 $g_2^{-1}g_1\in\operatorname{Stab}_G(x)$，恰好说明 $g_1$ 与 $g_2$ 在同一个左陪集里。也就是说，“所有把 $x$ 送到同一位置的群元”自动聚成一个稳定子陪集。

有限群情形于是得到最常用的计数式
\[
|G|=|\operatorname{Orb}_G(x)|\,|\operatorname{Stab}_G(x)|.
\]
以后点群分类中的极点计数、共轭类大小以及晶体等价位置数，都是这一个公式的不同外观。


\begin{exampleplain}[$D_3$ 作用在正三角形顶点上：轨道与等价原子]
\label{ex:d3-action-triangle-vertices}
令 $D_3$ 作用在正三角形顶点集 $X=\{1,2,3\}$ 上：$r$ 把 $1\to2\to3\to1$，$s$ 固定顶点 $1$ 而交换 $2\leftrightarrow3$。取顶点 $1$，其稳定子是所有保持 $1$ 不动的操作：
\[
\operatorname{Stab}_{D_3}(1)=\{e,s\}\cong C_2,
\]
几何上就是"过顶点 $1$ 的镜面反射加上恒等"，恰好是 $D_3$ 的一个 $C_2$ 子群。轨道是所有能从 $1$ 到达的顶点：
\[
\operatorname{Orb}_{D_3}(1)=\{1,\,r\cdot1=2,\,r^2\cdot1=3\}=\{1,2,3\}.
\]
轨道--稳定子定理给出
\[
\lvert\operatorname{Orb}_{D_3}(1)\rvert=\frac{|D_3|}{|\operatorname{Stab}_{D_3}(1)|}=\frac{6}{2}=3,
\]
即三个顶点。这里稳定子的陪集 $\{e\operatorname{Stab},\,r\operatorname{Stab},\,r^2\operatorname{Stab}\}$ 与三个顶点一一对应：陪集 $r\operatorname{Stab}$ 里的两个元素 $r$ 和 $rs=sr^2$ 都把顶点 $1$ 送到顶点 $2$，只是沿途经过的反射不同。

这个计数在分子物理中直接读作"等价原子个数"。$\mathrm{NH_3}$ 的三个 H 原子在 $C_{3v}\cong D_3$ 作用下彼此等价：旋转 $C_3$ 把一个 H 送到另一个，反射 $\sigma_v$ 交换两个 H。取任一 H 原子，其稳定子是过该原子的镜面反射（$C_s$，阶 $2$），轨道是全部三个 H（阶 $3$）。轨道--稳定子定理 $3=6/2$ 说明：三个 H 原子之所以等价，不是偶然的几何巧合，而是群作用把一个原子"复制"到了 $[G:\operatorname{Stab}]=3$ 个位置。后续用投影算符构造对称性适配线性组合时，每个等价原子轨道贡献一个基矢，等价原子个数因此直接进入对称适配线性组合的构造。
\end{exampleplain}

\begin{theorem}[Cayley 定理]
任意群 $G$ 都同构于某个置换群的子群。更具体地，$G$ 同构于 $\operatorname{Sym}(G)$ 的一个子群。
\label{thm:cayley}

\end{theorem}

\begin{proof}
让 $G$ 作用在自身上：
\[
L_g(x)=gx.
\]
由消去律，$L_g:G\to G$ 为双射，因此 $L_g\in\operatorname{Sym}(G)$。定义
\[
L:G\to\operatorname{Sym}(G),
\qquad
g\mapsto L_g.
\]
因为
\[
L_{gh}(x)=(gh)x=g(hx)=(L_g\circ L_h)(x),
\]
所以 $L$ 为同态。若 $L_g=\id(G)$，则特别有 $g=L_g(e)=e$，故 $\operatorname{Ker} L=\{e\}$，所以 $L$ 单射。于是
\[
G\cong L(G)\leqslant\operatorname{Sym}(G).
\]
\end{proof}

群表示本质上也是群作用，只是作用对象从一般集合 $X$ 换成了线性空间 $V$，并要求每个群元作用为可逆线性变换：
\[
\rho:G\to\operatorname{GL}(V).
\]
因此本节的核、忠实性、轨道与稳定子思想将在表示论和对称性物理中反复出现。

\begin{definitionplain}[群作用的分类]
\label{def:action-types}
设 $G$ 作用在 $X$ 上。
\begin{enumerate}[label=(\roman*),leftmargin=3em]
\item \textbf{传递作用}：$X$ 只有一个轨道，即对任意 $x,y\in X$ 存在 $g\in G$ 使 $g\cdot x=y$。
\item \textbf{自由作用}：每个稳定子平凡，即 $\operatorname{Stab}_G(x)=\{e\}$ 对所有 $x\in X$。等价地，只有 $e$ 有不动点。
\item \textbf{正则作用}：同时传递且自由。此时 $|X|=|G|$，且 $X$ 与 $G$ 自身作为 $G$-集等价。
\item \textbf{半正则作用}：所有稳定子相同（必为某个正规子群）。自由作用是半正则作用的特例。
\end{enumerate}
\end{definitionplain}

\begin{proposition}[传递作用的等价刻画]
\label{prop:transitive-action}
$G$ 在 $X$ 上传递作用，当且仅当 $X$ 作为 $G$-集同构于 $G/H$（左陪集空间），其中 $H$ 是某点的稳定子。具体地，取 $x_0\in X$，$H=\operatorname{Stab}_G(x_0)$，则
\[
X\cong G/H,\qquad g\cdot x_0\longmapsto gH.
\]
所有传递 $G$-集都由这种形式给出，且 $G/H_1\cong G/H_2$ 当且仅当 $H_1$ 与 $H_2$ 共轭。

\end{proposition}

\begin{proof}
轨道--稳定子定理已给出双射 $G/H\to\operatorname{Orb}(x_0)=X$（传递性使轨道等于全空间）。它是 $G$-等变的：$g'\cdot(gH\mapsto g\cdot x_0)=(g'gH\mapsto g'g\cdot x_0)$。

等价性：若 $\varphi:G/H_1\to G/H_2$ 是 $G$-等变双射，设 $\varphi(eH_1)=g_0H_2$。则对任意 $g$，
\[
\varphi(gH_1)=\varphi(g\cdot eH_1)=g\cdot\varphi(eH_1)=g g_0 H_2.
\]
$\varphi$ 良定义要求 $gH_1=eH_1\Rightarrow gg_0H_2=g_0H_2$，即 $H_1\subseteq g_0H_2g_0^{-1}$。由 $\varphi^{-1}$ 同理得反向包含，故 $H_1=g_0H_2g_0^{-1}$。
\end{proof}

\begin{theorem}[$p$-群作用的不动点定理]
\label{thm:pgroup-fixed-point}
设 $p$ 为素数，$P$ 为 $p$-群（$|P|=p^a$），$P$ 作用在有限集 $X$ 上。则不动点集
\[
X^P=\{x\in X\mid g\cdot x=x,\ \forall g\in P\}
\]
满足
\begin{equation}
|X^P|\equiv|X|\pmod{p}.
\label{eq:pgroup-fixed-point}
\end{equation}
特别地，若 $p\mid|X|$，则 $X^P\neq\varnothing$。

\end{theorem}

\begin{proof}
$X$ 分解为不相交的轨道之并。不动点恰是大小为 $1$ 的轨道。对非不动点 $x$，$|\operatorname{Orb}(x)|=[P:\operatorname{Stab}(x)]$ 整除 $|P|=p^a$ 且大于 $1$，故 $p\mid|\operatorname{Orb}(x)|$。因此
\[
|X|=\sum_{\text{不动点}}1+\sum_{\text{非不动点}}|\operatorname{Orb}(x_i)|\equiv|X^P|\pmod{p}.
\]
\end{proof}

\begin{corollary}[$p$-群有非平凡中心]
\label{cor:pgroup-center}
有限 $p$-群 $P$ 的中心 $Z(P)$ 非平凡，即 $|Z(P)|>1$。

\end{corollary}

\begin{proof}
$P$ 以共轭作用在自身上。不动点集恰是中心：
\[
P^P=\{x\in P\mid gxg^{-1}=x,\ \forall g\in P\}=Z(P).
\]
由 \cref{thm:pgroup-fixed-point}，$|Z(P)|\equiv|P|\equiv0\pmod{p}$。又 $e\in Z(P)$ 故 $|Z(P)|\ge1$，结合 $p\mid|Z(P)|$ 得 $|Z(P)|\ge p>1$。
\end{proof}

\begin{corollary}[$p^2$ 阶群是 Abel 群]
\label{cor:p2-abelian}
$|G|=p^2$ 的群是 Abel 群，因而同构于 $C_{p^2}$ 或 $C_p\times C_p$。

\end{corollary}

\begin{proof}
由 \cref{cor:pgroup-center}，$p\mid|Z(G)|$，故 $|Z(G)|\in\{p,p^2\}$。

若 $|Z(G)|=p^2$，则 $G=Z(G)$ 是 Abel 群。

若 $|Z(G)|=p$，则 $G/Z(G)$ 的阶为 $p$，故 $G/Z(G)\cong C_p$ 是循环群。设 $G/Z(G)=\langle gZ(G)\rangle$，则 $G$ 的每个元素形如 $g^k z$（$z\in Z(G)$）。对 $g^k z, g^\ell z'\in G$：
\[
(g^k z)(g^\ell z')=g^{k+\ell}zz'=(g^\ell z')(g^k z),
\]
因 $z,z'\in Z(G)$ 与一切元素交换。故 $G$ Abel，矛盾于 $|Z(G)|=p<p^2$。因此 $|Z(G)|=p$ 不可能，$G$ 必 Abel。

Abel $p^2$ 阶群的结构：取 $G$ 中最大阶元素 $g$。若 $\operatorname{ord}(g)=p^2$，则 $G\cong C_{p^2}$。若 $\operatorname{ord}(g)=p$，取与 $g$ 线性无关（在 $G$ 作为 $\mathbb F_p$-向量空间的意义下）的 $h$，则 $G\cong\langle g\rangle\times\langle h\rangle\cong C_p\times C_p$。
\end{proof}

\medskip
\noindent\emph{不动点定理的物理应用：晶场中的对称不动点。}
\label{ex:fixed-point-crystal}
在晶场理论中，晶格的对称群 $P$（一个 $p$-群的子群）作用在原子轨道集合 $X$ 上。不动点 $X^P$ 恰是在 $P$ 下不变的轨道——它们承载 $P$ 的平凡表示。

例如，$C_3$（$3$ 阶循环群）作用在 $p$ 轨道 $\{p_x,p_y,p_z\}$ 上（绕 $z$ 轴转 $120^\circ$）。$|X|=3$，$p=3\mid 3$，故不动点非空。事实上 $p_z$ 在 $C_3$ 下不变（$X^{C_3}=\{p_z\}$），而 $\{p_x,p_y\}$ 构成 $C_3$ 的二维轨道。不动点定理保证了 $p_z$ 的存在性，而不需要显式计算转动矩阵。

更一般地，在分子振动分析中，$p$-群不动点定理保证了某些振动模式在对称操作下必然不变——这些模式在光谱中表现为非简并的完全对称振动，是红外和 Raman 光谱中最容易识别的信号。
\par\medskip

群作用也可以在正三角形上完整算一次。令 $G=S_3$ 作用在顶点集合 $X=\{1,2,3\}$ 上。取顶点 $1$。它的轨道是全部三个顶点，
\[
\operatorname{Orb}_G(1)=\{1,2,3\},
\]
因为总能找到置换把 $1$ 送到任意顶点。稳定子则是所有保持 $1$ 不动的置换：
\[
\operatorname{Stab}_G(1)=\{e,(23)\},
\]
所以 $|\operatorname{Stab}_G(1)|=2$。轨道--稳定子定理给出
\[
|S_3|=|\operatorname{Orb}_G(1)|\,|\operatorname{Stab}_G(1)|=3\times2=6.
\]
如果换到共轭作用，完全相同的计数思想会把“顶点轨道”替换成“共轭类”，把“固定顶点的群元”替换成“与某群元交换的中心化子”。


群作用不仅能算单个轨道的大小，还能一次性算出\emph{轨道的总数}。这在物理和化学中极为常用：分子的不等价构型、晶格的不等价位点、配分函数中的对称因子，都可以用同一条公式处理。

\begin{theorem}[Burnside 引理]
设有限群 $G$ 作用在有限集 $X$ 上。轨道总数为
\[
\lvert X/G\rvert=\frac1{\lvert G\rvert}\sum_{g\in G}\lvert X^g\rvert,
\]
其中 $X^g=\{x\in X\mid g\cdot x=x\}$ 是 $g$ 的不动点集。
\label{thm:burnside}
\end{theorem}

这条公式的含义是：\emph{不等价构型数等于"平均不动点数"。} 每个群元各自数一下它保持不动的元素有多少个，再对所有群元取平均。

\par\medskip
\noindent\textbf{Burnside 引理的证明：双重计数}\quad
先构造二元组集合。对每个 $x\in X$，设 $r_x$ 为使 $x$ 不动的群元个数，即 $r_x=\lvert\operatorname{Stab}_G(x)\rvert$。考虑二元组集合
\[
\mathcal P=\{(g,x)\mid g\in G,\ x\in X,\ g\cdot x=x\}.
\]
这是所有"群元固定某元素"的配对。

随后按 $g$ 计数。固定 $g\in G$，满足 $g\cdot x=x$ 的 $x$ 恰有 $\lvert X^g\rvert$ 个（$g$ 的不动点集大小）。对 $g$ 求和：
\[
\lvert\mathcal P\rvert=\sum_{g\in G}\lvert X^g\rvert.
\]

接着按 $x$ 计数。固定 $x\in X$，使 $x$ 不动的 $g$ 恰有 $r_x=\lvert\operatorname{Stab}_G(x)\rvert$ 个（$x$ 的稳定子大小）。对 $x$ 求和：
\[
\lvert\mathcal P\rvert=\sum_{x\in X}r_x=\sum_{x\in X}\lvert\operatorname{Stab}_G(x)\rvert.
\]
因此
\[
\sum_{g\in G}\lvert X^g\rvert=\sum_{x\in X}\lvert\operatorname{Stab}_G(x)\rvert.
\]

再按轨道分组。对同一轨道中的元素，稳定子大小相同：若 $y=h\cdot x$，则 $\operatorname{Stab}(y)=h\operatorname{Stab}(x)h^{-1}$（共轭子群，同阶）。对轨道 $\operatorname{Orb}(x)$ 中每个元素，由轨道--稳定子定理
\[
\lvert\operatorname{Stab}_G(x)\rvert=\frac{\lvert G\rvert}{\lvert\operatorname{Orb}_G(x)\rvert}.
\]
故该轨道中所有元素的 $r_x$ 之和为
\[
\lvert\operatorname{Orb}_G(x)\rvert\cdot\frac{\lvert G\rvert}{\lvert\operatorname{Orb}_G(x)\rvert}=\lvert G\rvert.
\]
关键观察：轨道大小与稳定子大小恰好互消，每个轨道贡献恰好 $\lvert G\rvert$。

再汇总。设轨道总数为 $\lvert X/G\rvert$，每个轨道贡献 $\lvert G\rvert$，所以
\[
\sum_{g\in G}\lvert X^g\rvert=\lvert G\rvert\cdot\lvert X/G\rvert.
\]
两边除以 $\lvert G\rvert$ 即得 Burnside 公式。
\par\medskip

Burnside 引理的加权形式是 Pólya 计数定理。若每个位置可染 $c$ 种颜色，记 $c_i(g)$ 为置换 $g$ 中长度为 $i$ 的循环数，则染色轨道的生成函数为
\begin{equation*}
\frac{1}{|G|}\sum_{g\in G}\prod_{i}\left(x_1^i+x_2^i+\cdots+x_c^i\right)^{c_i(g)},
\end{equation*}
其中 $x_j$ 记录第 $j$ 种颜色；取 $x_j=1$ 就退化为 Burnside 引理。引入循环指标
\begin{equation*}
Z_G(t_1,\dots,t_n)=\frac{1}{|G|}\sum_{g\in G}t_1^{c_1(g)}t_2^{c_2(g)}\cdots t_n^{c_n(g)},
\end{equation*}
后，任意染色问题都可由代换 $t_i=x_1^i+\cdots+x_c^i$ 读出。另一方面，$\chi_{\mathrm{perm}}(g)=|X^g|$ 是置换表示的特征标，因此 Burnside 引理等价于
\begin{equation*}
|X/G|=\langle\chi_{\mathrm{perm}},\chi_{\mathrm{triv}}\rangle=\frac{1}{|G|}\sum_g\chi_{\mathrm{perm}}(g).
\end{equation*}
也就是说，轨道数正是置换表示中平凡表示的重数。

\medskip
\noindent\emph{用 Burnside 引理数正方形染色。}
正方形四个顶点用两种颜色着色，在 $D_4$ 的对称操作下有多少种不等价方案？

不加群作用时共有 $2^4=16$ 种着色。$D_4$ 有 $8$ 个元素，逐个数不动点：
\begin{itemize}
\item $e$：全部 $16$ 种都不动。
\item $r$（转 $90^\circ$）和 $r^3$（转 $270^\circ$）：四顶点必须同色，只有 $2$ 种。
\item $r^2$（转 $180^\circ$）：对角顶点同色，有 $2^2=4$ 种。
\item 两个关于对角线的反射：对角线上两点各自独立，另两点必须同色，有 $2^3=8$ 种。
\item 两个关于对边中点的反射：两对顶点各自同色，有 $2^2=4$ 种。
\end{itemize}
Burnside 引理给出
\[
\lvert X/G\rvert=\frac18(16+2+4+2+8+8+4+4)=\frac{48}{8}=6.
\]

同一计数也适用于统计物理中的对称等价构型，以及分子光谱中同位素取代位置的分类。
\par\medskip

\subsection{Cauchy 定理与 Sylow 定理}
\label{sec:cauchy-sylow}

Lagrange 定理告诉我们子群阶必须整除群阶，但\emph{反过来}不成立：$d\mid\lvert G\rvert$ 不保证存在 $d$ 阶子群。然而对素数幂阶子群，情况特别好——Sylow 定理保证了它们的存在性，并给出了共轭性和个数约束。这些结论是有限群结构理论的基石，也是后续分析分子对称群、晶体点群中"哪些子群必然存在"的工具。


Cauchy 定理是 Sylow 定理的最简单情形：只要素数 $p$ 整除群阶，群中就一定有 $p$ 阶元。它的证明是群作用计数技巧的典范。

\begin{theorem}[Cauchy 定理]
设 $G$ 为有限群，$p$ 为素数。若 $p\mid\lvert G\rvert$，则 $G$ 中存在阶为 $p$ 的元素，因而存在 $p$ 阶循环子群 $C_p\leqslant G$。
\label{thm:cauchy}
\end{theorem}

\par\medskip
\noindent\textbf{Cauchy 定理的证明：用 $p$ 元组上的 $\ZZ_p$ 作用}\quad
先构造 $p$ 元组集合。考虑
\[
S=\{(g_1,g_2,\dots,g_p)\mid g_1g_2\cdots g_p=e\}.
\]
自由选取 $g_1,\dots,g_{p-1}\in G$（共 $\lvert G\rvert^{p-1}$ 种选择），则 $g_p=(g_1\cdots g_{p-1})^{-1}$ 唯一确定且满足乘积条件。故 $\lvert S\rvert=\lvert G\rvert^{p-1}$。

随后$\ZZ_p$ 在 $S$ 上的循环移位作用。定义循环移位
\[
\tau\cdot(g_1,g_2,\dots,g_p)=(g_2,g_3,\dots,g_p,g_1).
\]
这个作用确实保持 $S$：若 $g_1g_2\cdots g_p=e$，则移位后
\[
g_2g_3\cdots g_p g_1=g_1^{-1}(g_1g_2\cdots g_p)g_1=g_1^{-1}\,e\,g_1=e,
\]
故移位后仍在 $S$ 中。循环移位 $p$ 次使每个分量回到原位，即 $\tau^p=\id()$，因此 $\ZZ_p=\langle\tau\rangle$ 作用在 $S$ 上。

接着轨道分析。$\ZZ_p$ 是素数阶群，其轨道大小（即 $\ZZ_p$ 的子群指数）只能是 $1$ 或 $p$。
\begin{itemize}
\item \emph{大小为 $1$ 的轨道}：循环移位不变的元组，即 $g_1=g_2=\cdots=g_p=g$，条件 $g^p=e$。设这样的元组共 $N$ 个（含 $g=e$ 给出的全 $e$ 元组）。
\item \emph{大小为 $p$ 的轨道}：其余元组，每个轨道含 $p$ 个不同元组。
\end{itemize}
故
\[
\lvert S\rvert = N + p\cdot(\text{非平凡轨道数}).
\]
因此 $N\equiv\lvert S\rvert\pmod{p}$。又 $\lvert S\rvert=\lvert G\rvert^{p-1}$ 且 $p\mid\lvert G\rvert$，故 $p\mid\lvert S\rvert$，从而 $p\mid N$。

最后，$N$ 是满足 $g^p=e$ 的元素数。单位元计入其中，且 $p\mid N$，故 $N\ge p$；于是存在非单位元 $g$ 满足 $g^p=e$。因为 $p$ 为素数，$g$ 的阶必为 $p$。
\par\medskip

Cauchy 定理把 Lagrange 定理对素数阶部分逆转过来；Sylow 第一定理进一步把结论提升到任意整除 $|G|$ 的素数幂。证明中“群作用加模计数”的结构也会在 Sylow 定理中再次出现。

另一个重要推论是：非平凡 $p$-群 $P$ 的中心必非平凡。类方程给出
\begin{equation*}
|P|=|Z(P)|+\sum_{a\notin Z(P)}[P:C_P(a)],
\end{equation*}
右端和式中每一项 $[P:C_P(a)]$ 都是 $p$ 的幂且 $>1$（因 $a\notin Z(P)$ 故 $C_P(a)\lneq P$），故每一项都被 $p$ 整除。于是
\begin{align*}
|Z(P)|&=|P|-\sum_{a\notin Z(P)}[P:C_P(a)]\\
&\equiv 0-0 \pmod{p}\\
&\equiv 0 \pmod{p},
\end{align*}
又因 $Z(P)$ 非空，所以 $|Z(P)|\ge p$。这为按群阶归纳研究 $p$-群提供了下降步骤。有限性不可省略；例如无限群 $\ZZ$ 除单位元外没有有限阶元素。

Cauchy 定理只保证 $p$ 阶元存在，但 Sylow 定理给出更强的结果：素数幂阶子群也必然存在，且共轭性和个数都有精确约束。

\begin{theorem}[Sylow 定理]
设 $G$ 为有限群，$\lvert G\rvert=p^a m$，$p$ 为素数，$p\nmid m$。
\begin{enumerate}[label=(\roman*),leftmargin=3em]
\item \textbf{存在性：}存在 $p^a$ 阶子群 $P\leqslant G$，称为 Sylow $p$-子群。
\item \textbf{共轭性：}任意两个 Sylow $p$-子群共轭：若 $P,Q$ 都是 Sylow $p$-子群，则存在 $g\in G$ 使 $Q=gPg^{-1}$。
\item \textbf{个数约束：}Sylow $p$-子群的个数 $n_p$ 满足
\[
n_p\equiv1\pmod{p},
\qquad
n_p\mid m.
\]
\end{enumerate}
\label{thm:sylow}
\end{theorem}

\par\medskip
\noindent\textbf{Sylow 定理的证明}\quad
(i) 存在性. 令 $\mathcal X$ 为 $G$ 中所有含 $p^a$ 个元素的子集，$G$ 以左乘作用在 $\mathcal X$ 上：$g\cdot S=gS$。子集总数
\[
|\mathcal X|=\binom{p^a m}{p^a}.
\]
这里需要同余 $\binom{p^a m}{p^a}\not\equiv0\pmod{p}$。把 $p^a m$ 个元素分成 $m$ 组、每组 $p^a$ 个。选取 $p^a$ 个元素时，若从第 $i$ 组选 $k_i$ 个（$\sum k_i=p^a$），贡献为 $\prod_i\binom{p^a}{k_i}$。模 $p$ 后，$\binom{p^a}{k_i}\equiv0\pmod{p}$ 除非 $k_i=0$ 或 $k_i=p^a$。因此只有从某一组全选的 $m$ 项留下：
\[
\binom{p^a m}{p^a}\equiv m\not\equiv0\pmod{p}.
\]

由轨道分解 $|\mathcal X|=\sum|\operatorname{Orb}_i|$，至少有一个轨道 $|\operatorname{Orb}(S)|$ 不被 $p$ 整除。由轨道--稳定子定理 $|\operatorname{Orb}(S)|=[G:\operatorname{Stab}(S)]$，故 $|\operatorname{Stab}(S)|=p^a m/|\operatorname{Orb}(S)|$。因 $p\nmid|\operatorname{Orb}(S)|$，$p^a\mid|\operatorname{Stab}(S)|$，即 $|\operatorname{Stab}(S)|\ge p^a$。

另一方面，取 $s\in S$，映射 $\operatorname{Stab}(S)\to S$，$h\mapsto hs$ 是单射（若 $h_1s=h_2s$ 则 $h_1=h_2$），故 $|\operatorname{Stab}(S)|\le|S|=p^a$。结合两个不等式，$|\operatorname{Stab}(S)|=p^a$，因此 $P=\operatorname{Stab}(S)$ 是 Sylow $p$-子群。

(ii) 共轭性. 设 $P,Q$ 都是 Sylow $p$-子群。让 $Q$ 以共轭作用在左陪集集合 $G/P=\{gP\mid g\in G\}$ 上：$q\cdot(gP)=qgP$。陪集数 $|G/P|=[G:P]=m$，$p\nmid m$。

由轨道分解 $m=\sum|\operatorname{Orb}(g_iP)|$。每个轨道大小 $|\operatorname{Orb}(g_iP)|=[Q:\operatorname{Stab}_Q(g_iP)]$ 整除 $|Q|=p^a$，故是 $p$ 的幂。但总和 $m$ 不被 $p$ 整除，因此至少有一个轨道大小为 $p^0=1$。设 $|\operatorname{Orb}(gP)|=1$，即对所有 $q\in Q$，$qgP=gP$，等价于 $g^{-1}qg\in P$，即 $g^{-1}Qg\subseteq P$。两边都是 $p^a$ 阶群，故 $g^{-1}Qg=P$，即 $Q=gPg^{-1}$。

(iii) 个数约束. 由共轭性，所有 Sylow $p$-子群彼此共轭，$n_p$ 是共轭作用的轨道数。让 $P$ 以共轭作用在 Sylow $p$-子群集合 $\{P_1,\dots,P_{n_p}\}$ 上。由轨道--稳定子定理，每个轨道大小整除 $|P|=p^a$，故是 $p$ 的幂。

这个作用的不动点只有 $P$ 自身。事实上，若 $Q$ 被 $P$ 固定，则 $P\le N_G(Q)$。此时 $P,Q$ 都是 $N_G(Q)$ 的 Sylow $p$-子群，而 $Q\mathrel{\trianglelefteq} N_G(Q)$；由已证的共轭性，只能有 $P=Q$。

因此 $P$ 的作用只有一个大小为 $1$ 的轨道（$\{P\}$），其余轨道大小都被 $p$ 整除。故
\[
n_p\equiv1\pmod{p}.
\]
又 $n_p=[G:N_G(P)]$（共轭轨道大小），而 $N_G(P)\ge P$，故
\[
n_p=[G:N_G(P)]\mid[G:P]=m.
\]
\par\medskip

最常用的推论是：$n_p=1$ 当且仅当 Sylow $p$-子群正规。实际判断有限群结构时，通常把 $n_p\equiv1\pmod p$ 与 $n_p\mid m$ 联立；若迫使某个 $n_p=1$，就立即得到一个非平凡正规子群。

\medskip
\noindent\emph{$D_3$ 的 Sylow 子群。}
$\lvert D_3\rvert=6=2\times3$。

对 $p=3$，有 $n_3\equiv1\pmod3$ 且 $n_3\mid2$，所以 $n_3=1$；因此
$\left\langle r\right\rangle=\{e,r,r^2\}$ 是唯一的 $3$ 阶子群并且正规。这与\cref{ex:d3-normal-subgroup-geometry} 的几何结论一致。

对 $p=2$，约束为 $n_2\equiv1\pmod2$ 且 $n_2\mid3$。群中确有
$\left\langle s\right\rangle,\left\langle sr\right\rangle,\left\langle sr^2\right\rangle$ 三个二阶子群，故 $n_2=3$；它们彼此共轭而都不正规。
\par\medskip

对阶为 $24$ 的点群，Sylow 约束立即给出 $n_2\in\{1,3\}$、$n_3\in\{1,4\}$；若 $n_3=4$，群对四个 Sylow $3$-子群的共轭作用给出到 $S_4$ 的同态。这样的计数常在进入具体表示计算前先确定可能的子群结构。

\medskip
\noindent\emph{$A_5$ 的单性证明——Sylow 计数的应用。}
\label{ex:a5-simplicity}
交错群 $A_5$ 有 $|A_5|=60=2^2\cdot3\cdot5$。用 Sylow 定理可以证明 $A_5$ 是单群（没有非平凡正规子群），这是 Galois 理论中"五次方程一般不可解"的核心事实。

先由 Sylow 定理核对各素数幂阶子群的数目：
\begin{itemize}
\item Sylow $5$-子群：$n_5\equiv1\pmod5$ 且 $n_5\mid12$。$A_5$ 中全部 $24$ 个 $5$-循环都是偶置换，每个 $C_5$ 含 $4$ 个非单位元，故 $n_5=6$。
\item Sylow $3$-子群：$n_3\equiv1\pmod3$ 且 $n_3\mid20$，故 $n_3\in\{1,4,10\}$。$A_5$ 中有 $20$ 个 $3$-循环（$(abc)$ 型），每个 $C_3$ 含 $2$ 个，故 $n_3=10$。
\item Sylow $2$-子群：$n_2\equiv1\pmod2$ 且 $n_2\mid15$，故 $n_2\in\{1,3,5,15\}$。$A_5$ 中 Klein 四元群 $V_4=\{e,(12)(34),(13)(24),(14)(23)\}$ 是 Sylow $2$-子群（$4$ 阶），共 $5$ 个（对应 $5$ 种把 $\{1,2,3,4,5\}$ 分成一对加单点的方式），故 $n_2=5$。
\end{itemize}

另一方面，正规子群必须是共轭类的并。$A_5$ 的共轭类大小为
\[
1,\quad 12,\quad 12,\quad 15,\quad 20,
\]
分别对应单位元、$20$ 个 $3$-循环、分成两类的 $24$ 个 $5$-循环，以及 $15$ 个双对换。

若 $N\mathrel{\trianglelefteq} A_5$，$|N|\mid60$ 且 $|N|=1+a\cdot20+b\cdot12+c\cdot12+d\cdot15$（$a,b,c,d\in\{0,1\}$，因每类要么全在要么全不在）。逐一检查 $|N|\in\{1,2,3,4,5,6,10,12,15,20,30,60\}$：
\begin{itemize}
\item $|N|=1$：$N=\{e\}$。
\item $|N|=60$：$N=A_5$。
\item $|N|=5$：需 $1+d\cdot15=5$ 或 $1+c\cdot12=5$ 等，均无整数解。
\item $|N|=10$：$1+a\cdot20=10$ 不行；$1+c\cdot12=10$ 不行；$1+d\cdot15=10$ 不行；$1+b\cdot12+c\cdot12=10$ 不行。
\item $|N|=15$：$1+d\cdot15=16\neq15$；$1+a\cdot20=15$ 不行；$1+b\cdot12=15$ 不行。
\item $|N|=20$：$1+a\cdot20=21\neq20$；$1+d\cdot15=20$ 不行；$1+b\cdot12=20$ 不行。
\item $|N|=30$：$1+a\cdot20+d\cdot15=30$ 即 $20a+15d=29$，无解；$1+b\cdot12+c\cdot12=30$ 即 $12(b+c)=29$，无解；$1+a\cdot20=30$ 不行。
\item $|N|=12$：$1+b\cdot12=13\neq12$；$1+c\cdot12=13\neq12$。
\item $|N|=6$：$1+a\cdot20=6$ 不行；$1+d\cdot15=6$ 不行。
\item $|N|=4$：$1+d\cdot15=4$ 不行。
\item $|N|=3$：$1+a\cdot20=3$ 不行。
\item $|N|=2$：$1+d\cdot15=2$ 不行。
\end{itemize}
因此 $A_5$ 没有非平凡正规子群，是单群。

$A_5$ 的单性意味着 $S_5$ 的合成列为
\[
S_5\rhd A_5\rhd\{e\},
\]
合成因子 $C_2$ 和 $A_5$。因 $A_5$ 非交换单群，$S_5$ 不可解。由 Galois 理论，五次方程一般不可根式求解。物理上，$A_5\cong I$（正二十面体转动群）的单性意味着正二十面体没有非平凡的连续对称性约化——它的全部对称性"不可拆分"。
\par\medskip

\subsection{合成列、Jordan--H\"older 定理与可解群}
\label{sec:composition-solvable}

Sylow 定理告诉我们素数幂阶子群必然存在，但不直接给出群的"建筑结构"。合成列和 Jordan--H\"older 定理填补这个空缺：任何有限群都可以被拆成一系列"不可再分"的片段（合成因子），而且这些片段在适当意义下是唯一的。可解群则是合成因子全是素数阶循环群的特殊情形，它直接联系到 Galois 理论中"方程何时可用根式求解"的经典判据。


\begin{definitionplain}[合成列]
群 $G$ 的合成列是严格递降的子群链
\[
G=G_0\rhd G_1\rhd\cdots\rhd G_n=\{e\},
\]
其中每个 $G_{i+1}$ 是 $G_i$ 的正规子群，且商群 $G_i/G_{i+1}$（称为合成因子）是\emph{单群}（没有非平凡正规子群）。长度 $n$ 称为合成列的长度。
\end{definitionplain}

\begin{theorem}[Jordan--H\"older]
设 $G$ 有两个合成列
\[
G=G_0\rhd G_1\rhd\cdots\rhd G_n=\{e\},
\qquad
G=H_0\rhd H_1\rhd\cdots\rhd H_m=\{e\}.
\]
则 $n=m$，且存在 $\{1,\dots,n\}$ 的置换 $\sigma$，使
\[
G_i/G_{i+1}\cong H_{\sigma(i)}/H_{\sigma(i)+1}.
\]
即合成因子的多重集（连同顺序的置换）与合成列的选择无关。
\label{thm:jordan-holder}
\end{theorem}

\par\medskip
\noindent\textbf{Jordan--H\"older 定理的证明}\quad
对 $|G|$ 归纳。若两条合成列的首个真子群相同，即 $G_1=H_1$，则对 $G_1$ 使用归纳假设即可。

设 $G_1\ne H_1$，并令 $N=G_1\cap H_1$。由于 $G/G_1$ 与 $G/H_1$ 都是单群，$G_1,H_1$ 是 $G$ 的极大正规子群。因此 $G_1H_1$ 只能是 $G$，而第二同构定理给出交叉同构
\[
G_1/N\cong G/H_1,
\qquad
H_1/N\cong G/G_1.
\]
取 $N$ 的任一合成列 $N=N_0\rhd\cdots\rhd N_k=\{e\}$。因为上式右端为单群，
\[
G_1\rhd N\rhd N_1\rhd\cdots\rhd\{e\},
\qquad
H_1\rhd N\rhd N_1\rhd\cdots\rhd\{e\}
\]
分别是 $G_1$ 与 $H_1$ 的合成列。对这两个较小的群应用归纳假设，可用它们替换原合成列的尾部。替换后，两边共有 $N$ 的全部合成因子；其余两项分别为
\[
\{G/G_1,\,G_1/N\},
\qquad
\{G/H_1,\,H_1/N\},
\]
而交叉同构表明这两个二元多重集相同。因此两条原合成列的合成因子仅相差排列，长度也相同。
\par\medskip

\medskip
\noindent\emph{$D_4$ 的合成列与合成因子。}
$D_4$（8 阶）有合成列
\[
D_4\rhd\left\langle r\right\rangle\rhd\left\langle r^2\right\rangle\rhd\{e\},
\]
其中 $r$ 为 $90^\circ$ 转动。商群：
\[
D_4/\left\langle r\right\rangle\cong C_2,
\qquad
\left\langle r\right\rangle/\left\langle r^2\right\rangle\cong C_2,
\qquad
\left\langle r^2\right\rangle/\{e\}\cong C_2.
\]
合成因子为 $\{C_2,C_2,C_2\}$。另一合成列 $D_4\rhd\left\langle r^2,s\right\rangle\rhd\left\langle s\right\rangle\rhd\{e\}$ 给出相同的 $\{C_2,C_2,C_2\}$，验证 Jordan--H\"older 定理。
\par\medskip

合成因子给出了群的"建筑砖块"，但什么样的群能从 Abel 砖块完全搭建出来？这就是可解群。

\begin{definitionplain}[可解群]
群 $G$ 是\emph{可解群}，若存在子群链
\[
G=G_0\rhd G_1\rhd\cdots\rhd G_n=\{e\},
\]
其中每个商群 $G_i/G_{i+1}$ 是 Abel 群。等价地，$G$ 的合成因子全是素数阶循环群 $C_p$。
\end{definitionplain}

\begin{proposition}[可解群的基本性质]
\begin{enumerate}[label=(\roman*),leftmargin=3em]
\item 可解群的子群和商群可解。
\item 可解群的可解扩张可解：若 $N\mathrel{\trianglelefteq} G$，$N$ 和 $G/N$ 都可解，则 $G$ 可解。
\item $p$-群可解（因 $p$-群有中心链，每步商为 $C_p$）。
\item $S_n$ 对 $n\ge5$ 不可解（因 $A_n$ 是单群且非 Abel）。
\end{enumerate}
\end{proposition}

\par\medskip
\noindent\textbf{可解群性质的证明}\quad
(i) 子群可解. 设 $H\le G$，$G$ 有 Abel 链 $G=G_0\rhd G_1\rhd\cdots\rhd G_n=\{e\}$。令 $H_i=H\cap G_i$。则
\[
H=H_0\ge H_1\ge\cdots\ge H_n=\{e\}.
\]
需验证每步正规且商 Abel。$H_{i+1}=H\cap G_{i+1}\mathrel{\trianglelefteq} H\cap G_i=H_i$：对 $h\in H_i$，$hG_{i+1}h^{-1}=G_{i+1}$（因 $G_{i+1}\mathrel{\trianglelefteq} G_i$ 且 $h\in G_i$），故 $hH_{i+1}h^{-1}=h(H\cap G_{i+1})h^{-1}=H\cap hG_{i+1}h^{-1}=H\cap G_{i+1}=H_{i+1}$。

商群 $H_i/H_{i+1}$ 嵌入 Abel 群 $G_i/G_{i+1}$：定义
\[
\psi:H_i\to G_i/G_{i+1},\qquad\psi(h)=hG_{i+1}.
\]
$\operatorname{Ker}\psi=H_i\cap G_{i+1}=H\cap G_i\cap G_{i+1}=H\cap G_{i+1}=H_{i+1}$。由第一同构定理 $H_i/H_{i+1}\hookrightarrow G_i/G_{i+1}$，Abel 群的子群 Abel，故 $H$ 可解。

(ii) 商群可解. 设 $N\mathrel{\trianglelefteq} G$，$G$ 有 Abel 链 $G_i$。令 $Q_i=G_iN/N$。由第二同构定理
\[
G_iN/N\cong G_i/(G_i\cap N),
\]
而 $G_i/(G_i\cap N)$ 是 $G_i/G_{i+1}$ 的商（因 $G_i\cap N\ge G_{i+1}$ 当 $N\mathrel{\trianglelefteq} G$ 时不一定成立，需用另一种方法）。更直接地，考虑投影 $\pi:G\to G/N$，令 $Q_i=\pi(G_i)=G_iN/N$。则
\[
G/N=Q_0\ge Q_1\ge\cdots\ge Q_n=\{eN\}.
\]
$Q_{i+1}\mathrel{\trianglelefteq} Q_i$ 因 $G_{i+1}\mathrel{\trianglelefteq} G_i$。商 $Q_i/Q_{i+1}$ 是 $G_i/G_{i+1}$ 的商群（由对应定理），Abel 群的商群 Abel。故 $G/N$ 可解。

(iii) 扩张可解. 设 $N\mathrel{\trianglelefteq} G$，$N$ 有 Abel 链 $N=N_0\rhd N_1\rhd\cdots\rhd N_k=\{e\}$，$G/N$ 有 Abel 链 $G/N=Q_0\rhd Q_1\rhd\cdots\rhd Q_m=\{eN\}$。令 $G_i=\pi^{-1}(Q_i)$（$\pi:G\to G/N$）。由对应定理 $G_i\mathrel{\trianglelefteq} G_{i-1}$ 且
\[
G_{i-1}/G_i\cong Q_{i-1}/Q_i
\]
是 Abel。拼接两条链：
\[
G=G_0\rhd G_1\rhd\cdots\rhd G_m=N\rhd N_1\rhd\cdots\rhd N_k=\{e\},
\]
每个商 Abel，故 $G$ 可解。

(iv) $p$-群可解. 设 $|G|=p^a$。由类方程
\[
|G|=|Z(G)|+\sum_{i}[G:C_G(g_i)],
\]
其中 $g_i$ 跑遍非中心共轭类代表。每个 $[G:C_G(g_i)]>1$ 是 $p$ 的幂，故 $p\mid|Z(G)|$，即中心非平凡。

取 $Z(G)$ 中 $p$ 阶子群 $C_p\le Z(G)$（由 Cauchy 定理，因 $p\mid|Z(G)|$）。因 $C_p\le Z(G)$，$C_p\mathrel{\trianglelefteq} G$。商群 $G/C_p$ 的阶为 $p^{a-1}$，仍是 $p$-群。由归纳假设 $G/C_p$ 可解，$C_p$ 亦可解，再由 (iii) 的扩张性质知 $G$ 可解。

归纳基础：$|G|=p^0=1$ 时 $G=\{e\}$ 可解（空链）。归纳步骤：$|G|=p^a$ 时，上述论证把问题降到 $p^{a-1}$ 阶，由归纳假设完成。

(v) $S_n$ 不可解（$n\ge5$）. $A_n$（$n\ge5$）是单群且非 Abel（\cref{ex:a5-simplicity} 对 $n=5$ 的证明可推广到一般 $n\ge5$）。若 $S_n$ 可解，则由 (i) 的子群遗传性，$A_n\le S_n$ 也应可解。但 $A_n$ 的合成列只能是
\[
A_n\rhd\{e\},
\]
因 $A_n$ 单。合成因子 $A_n$ 非 Abel，不满足可解群定义。矛盾，故 $S_n$ 不可解。
\par\medskip

\medskip
\noindent\emph{可解群与 Galois 理论的关联。}
方程 $f(x)=0$ 可用根式求解当且仅当其 Galois 群 $\operatorname{Gal}(f)$ 可解。

一般 $n$ 次方程的 Galois 群为 $S_n$。当 $n\le4$ 时，$S_n$ 可解（例如 $S_4$ 有 Abel 链 $S_4\rhd A_4\rhd V_4\rhd C_2\rhd\{e\}$），所以方程可由根式求解；当 $n\ge5$ 时，$S_n$ 不可解，这给出 Abel--Ruffini 定理的群论解释。

晶体点群也是可解群的重要来源：其合成因子为循环群，因此对称结构可以沿 Abel 商逐层分析。
\par\medskip


可解群要求合成因子是 Abel 群，但不要求这些 Abel 商如何嵌入原群。幂零群进一步加强：要求每一步商不仅 Abel，而且以\emph{中心}方式嵌入。这使得幂零群的结构比一般可解群更"刚性"。

\begin{definitionplain}[幂零群与中心列]
群 $G$ 是\emph{幂零群}，若存在子群链
\[
G=G_0\ge G_1\ge\cdots\ge G_n=\{e\},
\]
使得 $G_{i-1}/G_i\le Z(G/G_i)$ 对所有 $i$（即 $[G,G_{i-1}]\le G_i$）。这样的链称为\textbf{中心列}，最短长度称为\textbf{幂零类} $c(G)$。
\end{definitionplain}

条件 $G_{i-1}/G_i\le Z(G/G_i)$ 比 $G_{i-1}/G_i$ Abel 更强：它要求 $G_{i-1}/G_i$ 不仅交换，而且与整个 $G/G_i$ 交换。等价地，$[G,G_{i-1}]\le G_i$：换位子 $[g,h]$（$g\in G$，$h\in G_{i-1}$）落入 $G_i$。

\begin{definitionplain}[上中心列与下中心列]
$G$ 的\textbf{上中心列}定义为
\[
\{e\}=Z_0\le Z_1\le Z_2\le\cdots,\qquad Z_{i+1}/Z_i=Z(G/Z_i),
\]
即 $Z_1=Z(G)$，$Z_{i+1}$ 是 $G/Z_i$ 的中心在 $G$ 中的原像。\textbf{下中心列}定义为
\[
G=\gamma_1\ge\gamma_2\ge\cdots,\qquad\gamma_{i+1}=[G,\gamma_i],
\]
即 $\gamma_2=[G,G]=G'$（导群），$\gamma_3=[G,G']$，等等。
\end{definitionplain}

\begin{theorem}[幂零群的等价刻画]
\label{thm:nilpotent-equivalent}
对有限群 $G$，以下等价：
\begin{enumerate}[label=(\roman*),leftmargin=3em]
\item $G$ 幂零。
\item 上中心列终止于 $G$：存在 $c$ 使 $Z_c(G)=G$。
\item 下中心列终止于 $\{e\}$：存在 $c$ 使 $\gamma_{c+1}(G)=\{e\}$。
\item $G$ 是其 Sylow 子群的直积：$G=P_1\times P_2\times\cdots\times P_k$。
\end{enumerate}
满足上述条件时，幂零类 $c(G)$ 等于使 $Z_c=G$ 的最小 $c$，也等于使 $\gamma_{c+1}=\{e\}$ 的最小 $c$。
\end{theorem}

\begin{proof}
(i)$\Leftrightarrow$(ii). 若 $G$ 有中心列 $G=G_0\ge\cdots\ge G_n=\{e\}$，则 $G_i\le Z_{n-i}$（对 $n-i$ 归纳：$G_n=\{e\}=Z_0$；若 $G_i\le Z_{n-i}$，则 $[G,G_{i-1}]\le G_i\le Z_{n-i}$，故 $G_{i-1}/Z_{n-i}\le Z(G/Z_{n-i})$，即 $G_{i-1}\le Z_{n-i+1}$）。故 $G=G_0\le Z_n$，$Z_n=G$。逆方向类似。

(ii)$\Leftrightarrow$(iii). 对 $i$ 归纳可证 $\gamma_{i+1}\le Z_{c-i}$（$c$ 为上中心列长度）：$\gamma_1=G=Z_c$；若 $\gamma_i\le Z_{c-i+1}$，则 $\gamma_{i+1}=[G,\gamma_i]\le[G,Z_{c-i+1}]\le Z_{c-i}$（因 $Z_{c-i+1}/Z_{c-i}=Z(G/Z_{c-i})$ 蕴含 $[G,Z_{c-i+1}]\le Z_{c-i}$）。故 $\gamma_{c+1}\le Z_0=\{e\}$。逆方向类似。

(iii)$\Leftrightarrow$(iv). 这是有限群论的经典结果。\textbf{(iv)$\Rightarrow$(iii):} 直积的换位子 $[P_1\times\cdots\times P_k,P_1\times\cdots\times P_k]=[P_1,P_1]\times\cdots\times[P_k,P_k]$（不同 Sylow 分量交换），每个 $p$-群有下中心列终止于 $\{e\}$（因 $p$-群的中心非平凡，上中心列严格增长直到全群），故直积也如此。\textbf{(i)$\Rightarrow$(iv):} 幂零群的每个 Sylow 子群正规（因 Sylow $p$-子群在幂零群中是唯一的：若 $P,Q$ 都是 Sylow $p$-子群，$P\le Z_c$ 的某层但不在 $Z_{c-1}$，则 $Q$ 正规化 $P$，由 Sylow 共轭性 $P=Q$）。不同 Sylow 子群的交为 $\{e\}$（阶互质），元素彼此交换，故 $G$ 是直积。
\end{proof}

\begin{proposition}[幂零群的基本性质]
\label{prop:nilpotent-properties}
\begin{enumerate}[label=(\roman*),leftmargin=3em]
\item 幂零群的子群和商群幂零。
\item 幂零群的直积幂零，且 $c(G\times H)=\max(c(G),c(H))$。
\item $p$-群幂零，幂零类 $c\le\log_p|G|-1$。
\item 幂零群可解，但可解群不一定幂零。
\item 幂零群的每个 Sylow 子群正规（唯一 Sylow $p$-子群）。
\item 有限幂零群中，正规化子条件：$H\lneq N_G(H)$ 对所有真子群 $H$ 成立。
\end{enumerate}
\end{proposition}

\begin{proof}
(i) 子群：中心列与子群的交给出子群的中心列。商群：中心列取商给出商的中心列。

(ii) 直积的上中心列 $(Z_i(G\times H))=Z_i(G)\times Z_i(H)$，故 $c(G\times H)=\max(c(G),c(H))$。

(iii) $p$-群有非平凡中心（类方程），$Z_1=Z(G)\ne\{e\}$，$G/Z_1$ 仍是 $p$-群故有非平凡中心，$Z_2/Z_1=Z(G/Z_1)\ne\{e\}$，$Z_2$ 严格大于 $Z_1$。每步 $|Z_{i+1}/Z_i|\ge p$，故 $c\le\log_p|G|-1$。

(iv) 中心列的商 $G_{i-1}/G_i\le Z(G/G_i)$ 是 Abel 的，故中心列也是 Abel 链，幂零蕴含可解。反例：$S_3$ 可解（Abel 链 $S_3\rhd A_3\rhd\{e\}$）但不幂零（$Z(S_3)=\{e\}$，上中心列停滞在 $\{e\}$；或等价地，$S_3$ 的 Sylow $2$-子群不正规）。

(v) 由 \cref{thm:nilpotent-equivalent} (iv)，$G$ 是 Sylow 子群直积，每个因子唯一故正规。

(vi) 设 $H\lneq G$，取中心列 $G=Z_c\ge\cdots\ge Z_0=\{e\}$。设 $i$ 最大使 $Z_i\le H$（存在因 $Z_0=\{e\}\le H$，$Z_c=G\nleq H$）。则 $Z_{i+1}\nleq H$，但 $[G,Z_{i+1}]\le Z_i\le H$，故对 $z\in Z_{i+1}\setminus H$ 和 $h\in H$，$[z,h]=zhz^{-1}h^{-1}\in Z_i\le H$，即 $zhz^{-1}\in H$，$z\in N_G(H)$。故 $H\lneq\langle H,z\rangle\le N_G(H)$。
\end{proof}

\medskip
\noindent\emph{幂零群的实例与非实例。}
幂零:
\begin{itemize}
\item 所有 $p$-群：如 $D_4$（$2^3$ 阶），上中心列 $\{e\}\le\langle r^2\rangle\le\langle r\rangle\le D_4$，幂零类 $2$。
\item $D_4$ 幂零但 $D_3\cong S_3$ 不幂零：$D_3$ 的 Sylow $2$-子群（三个反射子群）不正规。
\item Abel 群幂零类 $0$ 或 $1$（$Z_1=G$）。
\item 直积 $C_4\times C_3\cong C_{12}$ 幂零（Sylow 分解 $C_4\times C_3$）。
\end{itemize}
不幂零:
\begin{itemize}
\item $S_3$：可解但不幂零（Sylow $2$-子群不正规）。
\item $S_4$：可解但不幂零（Sylow $3$-子群不正规）。
\item $A_4$：可解但不幂零（Sylow $2$-子群 $V_4$ 正规，但 Sylow $3$-子群不正规）。
\end{itemize}
\par\medskip

\medskip
\noindent\emph{Heisenberg 群：幂零类 $2$ 的典型例子。}
上三角整数矩阵群
\[
H_3(\ZZ)=\left\{\begin{pmatrix}1&a&c\\0&1&b\\0&0&1\end{pmatrix}:a,b,c\in\ZZ\right\}
\]
是幂零类 $2$ 的幂零群。计算：
\[
Z(H_3)=\left\{\begin{pmatrix}1&0&c\\0&1&0\\0&0&1\end{pmatrix}:c\in\ZZ\right\}\cong\ZZ,
\]
导群 $[H_3,H_3]=Z(H_3)$（换位子 $[x,y]$ 只在 $c$ 分量非零），故下中心列 $H_3\ge[H_3,H_3]=Z(H_3)\ge[H_3,Z(H_3)]=\{I\}$ 终止于第二步。上中心列 $\{I\}\le Z(H_3)\le H_3$ 也两步终止。

物理中，Heisenberg 群是量子力学位置--动量对易关系 $[x,p]=\ii\hbar$ 的群论化身：位置和动量的平移生成 $H_3$，中心分量是相空间面积（几何相位）。幂零性反映 $[x,p]$ 是中心元素（与一切交换），故"二次换位子"自动消失。
\par\medskip

可解群从内部子群链刻画群的"可分解性"，而自由群则从生成元和关系的角度给出群的"建筑图纸"。

\begin{definitionplain}[自由群]
集合 $S=\{s_1,\dots,s_n\}$ 上的自由群 $F(S)$ 是所有形式为
\[
s_{i_1}^{\epsilon_1}s_{i_2}^{\epsilon_2}\cdots s_{i_k}^{\epsilon_k},
\qquad\epsilon_j=\pm1,
\]
的\emph{既约字}（无相邻 $s_i s_i^{-1}$ 或 $s_i^{-1}s_i$ 可消去）构成的群，乘法为拼接后消去。$F(S)$ 满足泛性质：对任意群 $G$ 和任意映射 $\varphi:S\to G$，存在唯一同态 $\bar\varphi:F(S)\to G$ 扩展 $\varphi$。
\end{definitionplain}

\begin{definitionplain}[群的展示]
群 $G$ 的展示记为
\[
G=\langle S\mid R\rangle=\langle s_1,\dots,s_n\mid r_1,\dots,r_m\rangle,
\]
其中 $S$ 是生成元集，$R$ 是关系集（$F(S)$ 中的字）。$G\cong F(S)/\langle\!\langle R\rangle\!\rangle$，其中 $\langle\!\langle R\rangle\!\rangle$ 是 $R$ 生成的正规子群。
\end{definitionplain}

\medskip
\noindent\emph{常见群的展示。}
\begin{enumerate}[label=(\arabic*),leftmargin=2.5em]
\item 循环群：$C_n=\langle r\mid r^n=e\rangle$。
\item 二面体群：$D_n=\langle r,s\mid r^n=s^2=e,\,srs=r^{-1}\rangle$。
\item 对称群：$S_n=\langle\sigma_1,\dots,\sigma_{n-1}\mid\sigma_i^2=e,\,(\sigma_i\sigma_{i+1})^3=e,\,\sigma_i\sigma_j=\sigma_j\sigma_i\,(|i-j|\ge2)\rangle$（Coxeter 生成元）。
\item 四面体群：$T=\langle a,b\mid a^3=b^3=(ab)^2=e\rangle$。
\item 八面体群：$O=\langle a,b\mid a^4=b^3=(ab)^2=e\rangle$。
\end{enumerate}
\par\medskip

\par\medskip
\noindent\textbf{$D_n$ 展式的推导与验证}\quad
先从几何到代数。$D_n$ 由 $r$（$2\pi/n$ 转动）和 $s$（某反射）生成。几何关系：
\begin{itemize}
\item $r^n=e$：转 $n$ 次回原位；
\item $s^2=e$：反射幂等；
\item $srs=r^{-1}$：先转再反射等于反射再反转。
\end{itemize}
最后一个关系可从顶点作用直接验证。对 $v_k=\ee^{2\pi\ii k/n}$，取
$r(v_k)=v_{k+1}$ 与 $s(v_k)=v_{-k}$，便有
\[
srs(v_k)=s\bigl(r(v_{-k})\bigr)=s(v_{-k+1})=v_{k-1}=r^{-1}(v_k).
\]

随后化简任意字。设 $G=\langle r,s\mid r^n,s^2,srsr\rangle$。由关系 $srs=r^{-1}$，即 $sr=r^{-1}s$（右乘 $s$），可以把任意字中的 $s$ 逐步移到最右边：
\[
r^{a_1}s^{b_1}r^{a_2}s^{b_2}\cdots r^{a_k}s^{b_k}.
\]
因 $s^2=e$，每个 $b_i\in\{0,1\}$。因 $sr=r^{-1}s$，每次把 $s$ 穿过 $r$ 时翻转 $r$ 的幂次。最终化为
\[
r^k s^\epsilon,\qquad k\in\ZZ,\ \epsilon\in\{0,1\}.
\]
由 $r^n=e$，可取 $0\le k<n$。因此 $|G|\le 2n$（至多 $n$ 个 $r^k$ 和 $n$ 个 $r^k s$）。

正 $n$ 边形的几何对称群满足这些关系且有 $2n$ 个元素，所以泛性质给出满同态 $G\to D_n$。结合 $|G|\le2n$ 可知该同态为同构。

群的展示并不唯一。Tietze 变换给出两类保持群不变的改写：
\begin{itemize}
\item \emph{I 型}：添加一个可由现有关系导出的关系（或删除冗余关系）；
\item \emph{II 型}：添加一个新生成元 $a$ 及关系 $a=w$（$w$ 是旧生成元的字），或删除这样的生成元和关系。
\end{itemize}

从 $D_3=\langle r,s\mid r^3,s^2,srsr\rangle$ 出发，做 II 型变换：引入 $a=s$，$b=sr$。则 $r=sb=a^{-1}b=ab$（因 $a^2=e$）。新关系：
\begin{itemize}
\item $a^2=e$：由 $s^2=e$；
\item $b^2=(sr)^2=srsr=r^{-1}r=e$：由 $srs=r^{-1}$；
\item $(ab)^3=(s\cdot sr)^3=r^3=e$。
\end{itemize}
故 $D_3=\langle a,b\mid a^2,b^2,(ab)^3\rangle$，这是 Coxeter 形式。一般地
\[
D_n=\langle a,b\mid a^2,b^2,(ab)^n\rangle,
\]
直接反映 $D_n$ 由两个反射生成，其积有 $n$ 阶。

这个 Coxeter 形式也直接编码几何：$a,b$ 是夹角为 $\pi/n$ 的两个反射，其乘积 $ab$ 是转角 $2\pi/n$ 的旋转。

Coxeter 形式还直接给出 $D_n$ 的 Coxeter 图：两个顶点（代表两个反射）用标 $n$ 的边连接。$n=3$ 给 $D_3\cong S_3$，$n=4$ 给 $D_4$，$n=6$ 给 $D_6$。这些图在\chapref{chap:point-space-groups}的晶体点群分类中直接使用。
\par\medskip

因此从 $C_n$ 轴和 $\sigma_v$ 面满足的关系 $r^n=e$、$\sigma^2=e$、$\sigma r\sigma=r^{-1}$ 出发，就能重建 $C_{nv}$，不必先列出全部对称操作。

\subsection{直积与半直积}
\label{sec:products}


群还可以反过来由较小的群组装出来。最简单的情形是两个子系统彼此独立：一个群元写成 $(h,k)$，第一部分只按 $H$ 的规则相乘，第二部分只按 $K$ 的规则相乘，这就是直积。物理上“彼此独立”在代数里的表现是两个因子互相不改变对方。

但更常见的情况是一个子群会通过共轭作用改变另一个子群。例如二维刚体运动中，平移本身组成 Abel 群，但先转动再平移与先平移再转动并不一样：转动会把平移方向一起转过去。此时整体仍可由“平移部分 + 转动部分”描述，却不能使用普通直积；这正是半直积出现的原因。理解这一区别以后，空间群、二面体群以及许多物理对称群的乘法结构都会变得透明。


\begin{definitionplain}[群的直积]
设 $G,H$ 为群。笛卡尔积
\[
G\times H=\{(g,h)\mid g\in G,\ h\in H\}
\]
配备分量乘法
\[
(g_1,h_1)(g_2,h_2)
=(g_1g_2,h_1h_2)
\label{eq:direct-product}
\]
后构成群，称为 $G$ 与 $H$ 的直积。
\label{def:direct-product}
\end{definitionplain}

单位元为 $(e_G,e_H)$，逆元为
\[
(g,h)^{-1}=(g^{-1},h^{-1}).
\]
直积并不要求两个因子群本身是 Abel 群。

\begin{proposition}[直积中的标准正规子群]
在 $G\times H$ 中，
\[
G\times\{e_H\}\mathrel{\trianglelefteq} G\times H,
\qquad
\{e_G\}\times H\mathrel{\trianglelefteq} G\times H,
\]
且二者仅交于单位元，并逐元素交换。
\label{prop:direct-product-factors}

\end{proposition}

\begin{proof}
先看正规性。任取 $(g,h)\in G\times H$ 与 $(x,e_H)\in G\times\{e_H\}$，有
\[
(g,h)(x,e_H)(g,h)^{-1}
=(gxg^{-1},e_H)\in G\times\{e_H\}.
\]
所以 $G\times\{e_H\}\mathrel{\trianglelefteq} G\times H$；交换两个因子同理得到 $\{e_G\}\times H$ 正规。

若 $(x,e_H)=(e_G,y)$，比较两个分量立即得到 $x=e_G,y=e_H$，故二者交集只有 $(e_G,e_H)$。最后
\[
(x,e_H)(e_G,y)=(x,y)=(e_G,y)(x,e_H),
\]
所以两个标准因子中的元素逐元素交换。
\end{proof}

\begin{theorem}[内直积判据]
设 $N,M\mathrel{\trianglelefteq} G$，并满足
\[
N\cap M=\{e\},
\qquad
NM=G.
\]
则任意 $g\in G$ 都能唯一写成 $g=nm$，$n\in N,m\in M$，并且
\[
G\cong N\times M.
\label{eq:internal-direct-product}
\]
\label{thm:internal-direct-product}

\end{theorem}

\begin{proof}
先证 $N$ 与 $M$ 的元素逐元素交换。任取 $n\in N,m\in M$，交换子
\[
[n,m]=nmn^{-1}m^{-1}
\]
因 $N\mathrel{\trianglelefteq} G$ 而属于 $N$，又因 $M\mathrel{\trianglelefteq} G$ 而属于 $M$，故
\[
[n,m]\in N\cap M=\{e\}.
\]
于是 $nm=mn$。

存在性由 $NM=G$。若 $nm=n'm'$，则
\[
n'^{-1}n=m'm^{-1}.
\]
左边属于 $N$，右边属于 $M$，所以二者都属于 $N\cap M=\{e\}$，从而 $n=n'$、$m=m'$，分解唯一。

定义
\[
\Phi:N\times M\to G,
\qquad
\Phi(n,m)=nm.
\]
由逐元素交换，
\[
\Phi((n_1,m_1)(n_2,m_2))
=n_1n_2m_1m_2
=(n_1m_1)(n_2m_2),
\]
故 $\Phi$ 为同态；唯一分解说明它双射，因此为同构。
\end{proof}

\begin{corollary}[互素循环群的直积分解]
若 $\gcd(m,n)=1$，则
\[
C_m\times C_n\cong C_{mn}.
\label{eq:coprime-cyclic-product}
\]
特别地，
\[
C_6\cong C_2\times C_3.
\]
\label{cor:coprime-cyclic-product}

\end{corollary}

\begin{proof}
设 $a,b$ 分别生成 $C_m,C_n$。元素 $(a,b)$ 的阶是 $\operatorname{lcm}(m,n)=mn$，等于直积群的阶，因此它生成整个直积群。
\end{proof}

直积要求两个因子的群元互不"扭曲"对方。更一般地，可以允许 $H$ 通过自同构作用在 $N$ 上；这正是半直积的额外数据。

\begin{definitionplain}[外半直积]
设 $N,H$ 为群，并给定同态
\[
\varphi:H\longrightarrow\operatorname{Aut}(N).
\]
在集合 $N\times H$ 上定义乘法
\[
(n_1,h_1)(n_2,h_2)
=
\bigl(n_1\,\varphi(h_1)(n_2),\ h_1h_2\bigr).
\label{eq:semidirect-product}
\]
所得群称为由作用 $\varphi$ 定义的半直积，记为
\[
N\rtimes_{\varphi}H.
\]
\label{def:semidirect-product}
\end{definitionplain}

其单位元为 $(e_N,e_H)$。由\eqnrefs{eq:semidirect-product} 直接解方程可得
\[
(n,h)^{-1}
=
\bigl(\varphi(h^{-1})(n^{-1}),\ h^{-1}\bigr).
\label{eq:semidirect-inverse}
\]
当 $\varphi$ 为平凡同态，即 $\varphi(h)=\id(N)$ 对所有 $h$ 成立时，半直积退化为普通直积。


定义里把“所得群”三个字写得很短，但初学时不应跳过：必须验证这个乘法确实满足群公理。封闭性来自 $\varphi(h_1)(n_2)\in N$。结合律是唯一真正需要计算的一步。左边先算前两个因子：
\begin{align*}
&\bigl[(n_1,h_1)(n_2,h_2)\bigr](n_3,h_3)\\
&=\bigl(n_1\varphi(h_1)(n_2),h_1h_2\bigr)(n_3,h_3)\\
&=\Bigl(n_1\varphi(h_1)(n_2)\,
\varphi(h_1h_2)(n_3),\ h_1h_2h_3\Bigr).
\end{align*}
右边先算后两个因子：
\begin{align*}
&(n_1,h_1)\bigl[(n_2,h_2)(n_3,h_3)\bigr]\\
&=(n_1,h_1)\bigl(n_2\varphi(h_2)(n_3),h_2h_3\bigr)\\
&=\Bigl(n_1\varphi(h_1)\bigl(n_2\varphi(h_2)(n_3)\bigr),\ h_1h_2h_3\Bigr)\\
&=\Bigl(n_1\varphi(h_1)(n_2)\,
\varphi(h_1)\bigl(\varphi(h_2)(n_3)\bigr),\ h_1h_2h_3\Bigr).
\end{align*}
因为 $\varphi:H\to\operatorname{Aut}(N)$ 是群同态，
\[
\varphi(h_1h_2)=\varphi(h_1)\circ\varphi(h_2),
\]
两条计算的第一分量完全相同，故结合律成立。

单位元也可以直接代入检查：同态必把 $e_H$ 送到 $\operatorname{Aut}(N)$ 的单位元 $\id(N)$，所以
\[
(e_N,e_H)(n,h)=(n,h)=(n,h)(e_N,e_H).
\]
最后验证逆元而不是只“解方程”。令
\[
(n,h)^{-1}=\bigl(\varphi(h^{-1})(n^{-1}),h^{-1}\bigr).
\]
则
\begin{align*}
&(n,h)\bigl(\varphi(h^{-1})(n^{-1}),h^{-1}\bigr)\\
&=\Bigl(n\,\varphi(h)\bigl(\varphi(h^{-1})(n^{-1})\bigr),e_H\Bigr)\\
&=(nn^{-1},e_H)=(e_N,e_H),
\end{align*}
反向相乘同样得到单位元。到这里，外半直积的群结构才真正验证完毕。这个计算也显示了半直积和直积的唯一差别：第二个因子 $h_1$ 在乘进去以前会先通过 $\varphi(h_1)$ “扭动”后面的 $N$ 分量。

\begin{theorem}[内半直积判据]
设 $N\mathrel{\trianglelefteq} G$，$H\leqslant G$，并满足
\[
N\cap H=\{e\},
\qquad
NH=G.
\]
定义共轭作用
\[
\varphi:H\to\operatorname{Aut}(N),
\qquad
\varphi(h)(n)=hnh^{-1}.
\label{eq:conjugation-action}
\]
则
\[
G\cong N\rtimes_{\varphi}H.
\label{eq:internal-semidirect-product}
\]
\label{thm:internal-semidirect-product}

\end{theorem}

\begin{proof}
因 $N\mathrel{\trianglelefteq} G$，对每个 $h\in H$，共轭 $n\mapsto hnh^{-1}$ 把 $N$ 映到自身，且是自同构；又有
\[
\varphi(h_1h_2)=\varphi(h_1)\circ\varphi(h_2),
\]
故 $\varphi$ 为同态。

定义
\[
\Phi:N\rtimes_{\varphi}H\to G,
\qquad
\Phi(n,h)=nh.
\]
对两个元素，
\begin{align*}
\Phi(n_1,h_1)\Phi(n_2,h_2)
&=n_1h_1n_2h_2\\
&=n_1(h_1n_2h_1^{-1})h_1h_2\\
&=n_1\varphi(h_1)(n_2)h_1h_2\\
&=\Phi\bigl((n_1,h_1)(n_2,h_2)\bigr),
\end{align*}
所以 $\Phi$ 为同态。$NH=G$ 给出满射。若 $nh=e$，则 $n=h^{-1}\in N\cap H$，故 $n=h=e$，所以核平凡，$\Phi$ 单射。因此为同构。
\end{proof}

\medskip
\noindent\emph{$D_3\cong C_3\rtimes C_2$：反射如何"扭转"旋转。}
回到\cref{sec:group-axioms} 中正三角形 $D_3$ 的系统分析。取旋转子群和一条反射轴生成的子群：
\[
N=\left\langle r\right\rangle=\{e,r,r^2\}\cong C_3,
\qquad
H=\left\langle s\right\rangle=\{e,s\}\cong C_2.
\]
$N$ 是保持三角形手性的三个真转动，$H$ 由一条反射生成。由\cref{ex:d3-normal-subgroup-geometry}，$N\mathrel{\trianglelefteq} D_3$；又有 $N\cap H=\{e\}$ 以及 $|N||H|=|D_3|$，故 $NH=D_3$。

半直积中的作用由 $s$ 对 $r$ 的共轭决定。反射会翻转旋转方向，因此
\[
srs^{-1}=r^{-1}=r^2.
\]
于是 $C_2$ 对 $C_3$ 的作用为
\[
\varphi(s)(r^k)=r^{-k},
\qquad k=0,1,2.
\]
这是 $C_3$ 的非平凡自同构。内半直积判据\cref{thm:internal-semidirect-product} 因而给出
\[
D_3\cong C_3\rtimes_\varphi C_2,
\qquad
\varphi(s)=\text{"取逆"}.
\]

它不是直积，因为 $sr\neq rs$。半直积的乘法
\[
(r^{k_1},s^{\epsilon_1})(r^{k_2},s^{\epsilon_2})
=\bigl(r^{k_1+(-1)^{\epsilon_1}k_2},\,s^{\epsilon_1+\epsilon_2}\bigr)
\]
中的 $(-1)^{\epsilon_1}$ 正是反射翻转旋转方向的代数记录。对任意 $n\geq3$，同一论证给出
\[
D_n\cong C_n\rtimes C_2,
\qquad
\varphi(s)(r^k)=r^{-k}.
\]
旋转部分 $C_n$ 是正规子群（反射共轭旋转等于旋转的逆），而反射部分 $C_2$ 通过"反转旋转方向"非平凡地作用。这一结构是所有二面体群的代数骨架：\emph{旋转提供循环因子，反射提供二阶因子，半直积的作用记录"反射如何翻转旋转方向"}。
\par\medskip

直积与半直积都把复杂群拆成较小因子，但信息量不同。直积只需要两个因子群；半直积除了 $N$ 与 $H$，还必须给出作用
\[
\varphi:H\to\operatorname{Aut}(N).
\]
同一对因子群配上不同的 $\varphi$，可能得到不同构的群。因此“因子是什么”与“因子之间如何作用”是两层独立信息。

\subsection{习题与解答}
\label{sec:chapter-summary-exercises}

先用几道基础题把群公理、阶与子群判据真正算熟。它们的共同目标是让后面关于陪集、正规子群和作用的论证不再依赖直觉。

\begin{exercise}[唯一的三阶群]
证明：任意三阶群都同构于 $C_3$。
\end{exercise}

\emph{解答.} 取任意非单位元 $a\in G$。由 Lagrange 定理，$\operatorname{ord}(a)$ 必整除 $|G|=3$；又因 $a\neq e$，故 $\operatorname{ord}(a)\neq1$，只能有
\[
\operatorname{ord}(a)=3.
\]
因此循环子群 $\left\langle a\right\rangle$ 已有三个元素，阶与 $G$ 相同，所以 $G=\left\langle a\right\rangle\cong C_3$。


\begin{exercise}[子群的交]
不用直接引用\cref{prop:subgroup-intersection}，从一步子群判据重新证明两个子群的交仍为子群。并说明两个子群的并一般为什么不是子群。
\end{exercise}

\emph{解答.} 设 $H,K\leqslant G$。因为 $e\in H$ 且 $e\in K$，故 $H\cap K$ 非空。若 $a,b\in H\cap K$，则
\[
ab^{-1}\in H,
\qquad
ab^{-1}\in K,
\]
所以 $ab^{-1}\in H\cap K$。由一步子群判据，$H\cap K\leqslant G$。

并集一般不封闭。例如在加法群 $\ZZ$ 中，$2\ZZ$ 与 $3\ZZ$ 都是子群，但
\[
2\in2\ZZ,
\qquad3\in3\ZZ,
\qquad2+3=5\notin2\ZZ\cup3\ZZ.
\]
事实上，$H\cup K$ 是子群当且仅当 $H\subseteq K$ 或 $K\subseteq H$。


\begin{exercise}[四阶群]
证明：任意四阶群都是 Abel 群，并且恰有两种同构类型：$C_4$ 与 $C_2\times C_2$。
\end{exercise}

\emph{解答.} 设 $|G|=4$。若存在 $a\in G$ 满足 $\operatorname{ord}(a)=4$，则 $G=\left\langle a\right\rangle\cong C_4$。

若不存在四阶元，则由 Lagrange 定理，每个非单位元的阶只能为 $2$。任取 $a,b\in G$，若其中一个为 $e$，由单位元性质二者交换；否则 $a^{-1}=a$、$b^{-1}=b$，且 $ab$ 也满足 $(ab)^{-1}=ab$。另一方面
\[
(ab)^{-1}=b^{-1}a^{-1}=ba,
\]
故 $ab=ba$，所以 $G$ 为 Abel 群。取两个不同非单位元 $a,b$，则 $ab$ 既不是 $e$、$a$、也不是 $b$，因此
\[
G=\{e,a,b,ab\},
\qquad a^2=b^2=(ab)^2=e,
\]
故 $G\cong C_2\times C_2$。两种同构类型恰为 $C_4$ 与 Klein 四元群。


\begin{exercise}[全为二阶元的群]
设 $G$ 中每个非单位元都满足 $g^2=e$。证明 $G$ 为 Abel 群。
\end{exercise}

\emph{解答.} 对任意 $a,b\in G$，若 $ab=e$，则 $b=a^{-1}=a$，当然有 $ab=ba$。若 $ab\neq e$，由条件 $(ab)^2=e$，所以 $(ab)^{-1}=ab$。又
\[
(ab)^{-1}=b^{-1}a^{-1}=ba
\]
且 $a^{-1}=a,b^{-1}=b$，于是 $ab=ba$。任意两元素交换，故 $G$ 为 Abel 群。


\begin{exercise}[陪集不是子群]
设 $H\leqslant G$，$g\notin H$。证明 $gH$ 不含 $H$ 的任何元素，并说明 $gH$ 不可能是子群。
\end{exercise}

\emph{解答.} 若存在 $x\in H\cap gH$，则可写 $x=gh$，其中 $h\in H$。于是
\[
g=xh^{-1}\in H,
\]
与 $g\notin H$ 矛盾，所以 $H\cap gH=\varnothing$。特别地 $e\in H$，故 $e\notin gH$。任何子群都必须包含单位元，因此 $gH$ 不可能是子群。


\begin{exercise}[循环群的子群]
设 $G=C_n=\left\langle a\right\rangle$。证明：对每个正因子 $d\mid n$，$G$ 恰有一个 $d$ 阶子群，并写出其生成元。
\end{exercise}

\emph{解答.} 写 $n=dq$。元素 $a^q$ 的阶为
\[
\operatorname{ord}(a^q)=\frac{n}{\gcd(n,q)}=\frac{n}{q}=d,
\]
所以 $H_d=\left\langle a^q\right\rangle$ 是一个 $d$ 阶子群。

再证唯一性。设 $H\leqslant C_n$ 且 $|H|=d$。取最小正整数 $k$ 使 $a^k\in H$。对任意 $a^m\in H$，作带余除法 $m=q'k+r$，$0\le r<k$，则
\[
a^r=a^m(a^k)^{-q'}\in H.
\]
由 $k$ 的最小性只能有 $r=0$，故 $H=\left\langle a^k\right\rangle$。于是 $\operatorname{ord}(a^k)=d$，即
\[
\frac{n}{\gcd(n,k)}=d.
\]
最小正生成指数必为 $k=n/d=q$，所以 $H=H_d$。


下面几题转向共轭、正规子群与同态结构，重点是练习如何从已经证明的一般定理出发组织证明，而不是逐项检查群公理。

\begin{exercise}[$S_3$ 的子群与正规子群]
列出 $S_3$ 的全部子群，判断哪些正规。再直接构造 $S_3/A_3\cong C_2$ 的同构。
\end{exercise}

\emph{解答.} $S_3$ 的子群为
\[
\{e\},\quad
\left\langle (12)\right\rangle,\quad
\left\langle (13)\right\rangle,\quad
\left\langle (23)\right\rangle,\quad
A_3=\left\langle (123)\right\rangle,\quad
S_3.
\]
三个二阶子群彼此由共轭变换交换，例如
\[
(123)(12)(123)^{-1}=(23),
\]
所以它们都不正规。$A_3$ 的指数为 $2$，故正规；$\{e\}$ 与 $S_3$ 当然正规。

符号同态
\[
\operatorname{sgn}:S_3\to\{\pm1\}\cong C_2
\]
为满同态，核为 $A_3$。由第一同构定理，
\[
S_3/A_3\cong C_2.
\]


\begin{exercise}[共轭保持群元阶]
证明：若 $x\sim y$，则 $\operatorname{ord}(x)=\operatorname{ord}(y)$。进一步说明为什么同一个共轭类中的群元具有相同的循环型不变量。
\end{exercise}

\emph{解答.} 若 $y=gxg^{-1}$，则对任意整数 $n\ge1$，
\[
y^n=(gxg^{-1})^n=gx^ng^{-1}.
\]
因此 $y^n=e$ 当且仅当 $x^n=e$，两者满足单位元的最小正指数相同，所以 $\operatorname{ord}(y)=\operatorname{ord}(x)$。

对置换群，若 $x=(a_1\cdots a_k)$，则
\[
g x g^{-1}=(g(a_1)\cdots g(a_k)),
\]
共轭只重新命名标签，不改变循环长度，因此整个循环类型保持不变。


\begin{exercise}[逆共轭类]
若 $C$ 是 $G$ 的一个共轭类，定义
\[
C^{-1}=\{x^{-1}\mid x\in C\}.
\]
证明 $C^{-1}$ 也是一个共轭类。
\end{exercise}

\emph{解答.} 若 $C=\operatorname{Cl}(x)$，则任意 $y\in C$ 可写为 $y=gxg^{-1}$，从而
\[
y^{-1}=gx^{-1}g^{-1}.
\]
于是
\[
C^{-1}=\{gx^{-1}g^{-1}\mid g\in G\}=\operatorname{Cl}(x^{-1}),
\]
所以它仍是一个共轭类。


\begin{exercise}[唯一二阶元在中心中]
若群 $G$ 恰有一个二阶元素 $h$，证明
\[
h\in Z(G).
\]
\end{exercise}

\emph{解答.} 对任意 $g\in G$，共轭元 $ghg^{-1}$ 与 $h$ 有相同阶，因此也是二阶元。由二阶元唯一性，
\[
ghg^{-1}=h.
\]
右乘 $g$ 得 $gh=hg$。由于 $g$ 任意，$h\in Z(G)$。


\begin{exercise}[核的逆像描述]
设 $\varphi:G\to H$ 为同态，$K\leqslant H$。证明 $\varphi^{-1}(K)\leqslant G$；若 $K\mathrel{\trianglelefteq} H$，进一步证明 $\varphi^{-1}(K)\mathrel{\trianglelefteq} G$。
\end{exercise}

\emph{解答.} 因 $e_H\in K$ 且 $\varphi(e_G)=e_H$，所以 $\varphi^{-1}(K)$ 非空。若 $a,b\in\varphi^{-1}(K)$，则
\[
\varphi(ab^{-1})=\varphi(a)\varphi(b)^{-1}\in K,
\]
故 $ab^{-1}\in\varphi^{-1}(K)$，由一步子群判据它是 $G$ 的子群。

若进一步 $K\mathrel{\trianglelefteq} H$，取 $g\in G$、$x\in\varphi^{-1}(K)$，有
\[
\varphi(gxg^{-1})
=\varphi(g)\varphi(x)\varphi(g)^{-1}\in K,
\]
故 $gxg^{-1}\in\varphi^{-1}(K)$，所以 $\varphi^{-1}(K)\mathrel{\trianglelefteq} G$。


\begin{exercise}[第一同构定理的计算例]
定义
\[
\varphi:\ZZ\to C_n,
\qquad
k\mapsto a^k,
\]
其中 $C_n=\left\langle a\right\rangle$。求 $\operatorname{Ker}\varphi$ 与 $\operatorname{Im}\varphi$，并由第一同构定理证明
\[
\ZZ/n\ZZ\cong C_n.
\]
\end{exercise}

\emph{解答.} 因 $a$ 生成 $C_n$，$\varphi$ 为满射，所以 $\operatorname{Im}\varphi=C_n$。又
\[
\varphi(k)=e
\Longleftrightarrow a^k=e
\Longleftrightarrow n\mid k,
\]
故
\[
\operatorname{Ker}\varphi=n\ZZ.
\]
第一同构定理立即给出
\[
\ZZ/n\ZZ
=\ZZ/\operatorname{Ker}\varphi
\cong\operatorname{Im}\varphi=C_n.
\]


最后把群作用、Cayley 定理以及直积和半直积串起来。这些题直接为\chapref{chap:representation-theory}的表示作用语言做准备。

\begin{exercise}[轨道是等价类]
在 $X$ 上定义关系
\[
x\sim y
\quad\Longleftrightarrow\quad
\exists g\in G:\ y=g\cdot x.
\]
证明这是等价关系，并证明其等价类正是 $G$ 的轨道。
\end{exercise}

\emph{解答.} 自反性来自 $e\cdot x=x$。若 $y=g\cdot x$，则
\[
x=g^{-1}\cdot y,
\]
故对称性成立。若 $y=g\cdot x$ 且 $z=h\cdot y$，则
\[
z=h\cdot(g\cdot x)=(hg)\cdot x,
\]
故传递性成立。因此 $\sim$ 是等价关系。由定义，与 $x$ 等价的全部点正是
\[
\{g\cdot x\mid g\in G\}=\operatorname{Orb}_G(x),
\]
即轨道。


\begin{exercise}[共轭类作为轨道]
令 $G$ 在自身上按共轭作用
\[
g\cdot x=gxg^{-1}.
\]
证明 $x$ 的轨道为 $\operatorname{Cl}(x)$，稳定子为 $C_G(x)$，并用轨道--稳定子定理重新推出\eqnrefs{eq:class-size}。
\end{exercise}

\emph{解答.} 共轭作用下
\[
\operatorname{Orb}_G(x)=\{gxg^{-1}\mid g\in G\}=\operatorname{Cl}(x).
\]
稳定子为
\[
\operatorname{Stab}_G(x)=\{g\in G\mid gxg^{-1}=x\}
=\{g\in G\mid gx=xg\}=C_G(x).
\]
由轨道--稳定子定理，
\[
|\operatorname{Cl}(x)|=[G:C_G(x)]
=\frac{|G|}{|C_G(x)|},
\]
即共轭类大小公式。


\begin{exercise}[Cayley 定理的右正则版本]
定义右作用 $R_g(x)=xg^{-1}$。证明 $g\mapsto R_g$ 仍给出 $G$ 到 $\operatorname{Sym}(G)$ 的单射同态。解释为什么若直接取 $x\mapsto xg$，乘法次序会反转。
\end{exercise}

\emph{解答.} 定义 $R_g(x)=xg^{-1}$。则
\begin{align*}
(R_g\circ R_h)(x)
&=R_g(xh^{-1})\\
&=xh^{-1}g^{-1}\\
&=x(gh)^{-1}=R_{gh}(x),
\end{align*}
所以 $g\mapsto R_g$ 是同态。若 $R_g=R_h$，作用于 $e$ 得 $g^{-1}=h^{-1}$，故 $g=h$，因此它是单射。

若直接定义 $T_g(x)=xg$，则
\[
(T_g\circ T_h)(x)=xhg=T_{hg}(x),
\]
乘法顺序反转，所以 $g\mapsto T_g$ 是反同态而不是同态。


\begin{exercise}[直积的商群]
证明
\[
(G\times H)/(G\times\{e_H\})\cong H,
\qquad
(G\times H)/(\{e_G\}\times H)\cong G.
\]
要求明确写出满同态及其核。
\end{exercise}

\emph{解答.} 定义投影
\[
\pi_H:G\times H\to H,
\qquad
\pi_H(g,h)=h.
\]
它是满同态，且
\[
\operatorname{Ker}\pi_H=G\times\{e_H\}.
\]
第一同构定理给出
\[
(G\times H)/(G\times\{e_H\})\cong H.
\]
同理取 $\pi_G(g,h)=g$，其核为 $\{e_G\}\times H$，得到第二个同构。


\begin{exercise}[半直积中的正规因子]
在 $N\rtimes_\varphi H$ 中证明
\[
N\times\{e_H\}\mathrel{\trianglelefteq} N\rtimes_\varphi H,
\]
并说明 $\{e_N\}\times H$ 何时正规。由此给出半直积退化为直积的一个充要条件。
\end{exercise}

\emph{解答.} 把 $N$ 识别为 $N\times\{e_H\}$。对任意 $(n,h)$ 与 $(m,e_H)$，利用半直积乘法与逆元公式可得
\[
(n,h)(m,e_H)(n,h)^{-1}
=\bigl(n\varphi(h)(m)n^{-1},e_H\bigr)\in N\times\{e_H\},
\]
故该因子正规。

再看 $H_0=\{e_N\}\times H$。取 $(n,e_H)$ 与 $(e_N,h)$，
\[
(n,e_H)(e_N,h)(n,e_H)^{-1}
=\bigl(n\varphi(h)(n^{-1}),h\bigr).
\]
它对所有 $n,h$ 都属于 $H_0$ 当且仅当
\[
\varphi(h)(n)=n,
\qquad \forall n\in N,
\]
也就是 $\varphi$ 为平凡作用。因此 $H_0$ 正规当且仅当半直积退化为直积 $N\times H$。


\begin{exercise}[$D_4$ 的半直积结构]
设
\[
D_4=\left\langle r,s\mid r^4=s^2=e,\ srs^{-1}=r^{-1}\right\rangle.
\]
证明
\[
D_4\cong C_4\rtimes C_2,
\]
并明确写出 $C_2\to\operatorname{Aut}(C_4)$ 的非平凡作用。
\end{exercise}

\emph{解答.} 令 $N=\left\langle r\right\rangle\cong C_4$、$H=\left\langle s\right\rangle\cong C_2$。关系
\[
srs^{-1}=r^{-1}
\]
说明 $N\mathrel{\trianglelefteq} D_4$，而 $N\cap H=\{e\}$、$NH=D_4$。因此内半直积判据给出
\[
D_4\cong C_4\rtimes_\varphi C_2.
\]
作用由生成元 $s$ 决定：
\[
\varphi(s)(r^k)=r^{-k}.
\]
它是 $C_4$ 的非平凡二阶自同构。由于该作用不是恒等作用，$D_4$ 不是 $C_4\times C_2$。
